\documentclass[11pt,a4paper]{article}
\usepackage[T1]{fontenc}
\usepackage[utf8]{inputenc}
\usepackage[margin=21mm,headheight=15pt]{geometry}
\usepackage{lmodern,tabularx,booktabs,xcolor,enumitem,fancyhdr,hyperref}
\definecolor{navy}{HTML}{0B2756}
\definecolor{blue}{HTML}{0984FD}
\definecolor{muted}{HTML}{526473}
\hypersetup{colorlinks=true,urlcolor=blue,pdftitle={Ringkasan Serah Terima Proyek Magang SAPA},pdfauthor={Muhammad Naufal Tiftazani}}
\renewcommand{\familydefault}{\sfdefault}
\setlength{\parindent}{0pt}
\setlength{\parskip}{6pt}
\setlist[itemize]{leftmargin=16pt,itemsep=3pt,topsep=3pt}
\renewcommand{\arraystretch}{1.18}
\pagestyle{fancy}
\fancyhf{}
\fancyhead[L]{\small\color{navy}\textbf{SAPA | Proyek Magang}}
\fancyhead[R]{\small\color{muted}Pendamping presentasi}
\fancyfoot[L]{\footnotesize\color{muted}Draf review pembimbing | 26 Oktober 2026}
\fancyfoot[R]{\footnotesize\thepage/3}
\renewcommand{\headrulewidth}{0.4pt}
\newcommand{\topic}[1]{\vspace{5pt}{\large\color{navy}\bfseries #1}\par\vspace{2pt}}
\newcolumntype{Y}{>{\raggedright\arraybackslash}X}

\begin{document}
{\fontsize{28}{31}\selectfont\bfseries\color{navy}SAPA}\par
{\large Sistem Asisten Pelayanan \& Administrasi}\par
{\color{muted}Ringkasan serah terima hasil proyek magang untuk pembimbing lapangan}

\vspace{4pt}
\begin{tabularx}{\linewidth}{@{}p{36mm}Y@{}}
\toprule
Mahasiswa / NIM & Muhammad Naufal Tiftazani / 230705117 \\
Instansi / unit & Diskominfotik Kota Banda Aceh / E-Government \\
Fokus laporan & API Gateway dan sistem manajemen dokumen \\
Pembimbing & \textit{Diisi sebelum penyerahan} \\
Tanggal / versi paket & \textit{Diisi saat penyerahan / setelah versi kode dibekukan} \\
\bottomrule
\end{tabularx}

\topic{1. Tujuan dan manfaat}
SAPA dikembangkan untuk membantu layanan informasi yang bersumber dari
dokumen resmi dinas. Presentasi yang telah disampaikan menunjukkan pengalaman
penelusuran chat yang berakhir pada pesan petugas tidak tersedia dan formulir.
Hal itu menjadi motivasi proyek, bukan hasil survei seluruh layanan.

Manfaat yang dituju adalah memudahkan pemeliharaan basis pengetahuan,
membantu menjawab pertanyaan berulang, menampilkan sumber jawaban, serta
menyediakan jalur bantuan petugas ketika diperlukan.

\topic{2. Alur sistem yang diserahkan}
\begin{itemize}
\item Admin mengunggah dokumen, memeriksa hasil ekstraksi, dan memperbaiki teks.
\item Teks yang disetujui diolah menjadi basis pengetahuan untuk pencarian.
\item Pertanyaan diproses menggunakan konteks dokumen dan model AI.
\item Tiket membantu tindak lanjut yang memerlukan penanganan petugas.
\end{itemize}

\begin{tabularx}{\linewidth}{@{}p{38mm}Y@{}}
\toprule
Komponen & Tanggung jawab \\
\midrule
Admin frontend & Dashboard dokumen, koreksi, pengaturan, percakapan dan monitoring. \\
Laravel Admin & Autentikasi, API Gateway, data dokumen/tiket dan integrasi kanal. \\
Python AI & Ekstraksi, embedding, pencarian dokumen, respons AI dan worker. \\
\bottomrule
\end{tabularx}

\vfill
{\small\color{muted}Identitas mengikuti laporan yang tersedia. Kontribusi pribadi
perlu dijelaskan sesuai pekerjaan/logbook aktual; cakupan sistem tidak
menyatakan seluruh modul dikerjakan oleh satu orang.}

\newpage
\topic{3. Hubungan dengan presentasi sebelumnya}
Dokumen ini mendampingi \textbf{PPT DISKOMINFOTIK.pptx} yang berisi 18 slide.
Alurnya tetap sama: observasi, kebutuhan, ide layanan, arsitektur, lalu demo.

\begin{tabularx}{\linewidth}{@{}p{24mm}Y@{}}
\toprule
Slide & Penjelasan pendamping \\
\midrule
1--8 & Masalah yang diamati serta kebutuhan warga, petugas dan dinas. \\
9--11 & Gagasan hubungan pertanyaan, dokumen resmi dan bantuan petugas. \\
12 & Unggah, ekstraksi, koreksi manual dan pengindeksan dokumen. \\
13--17 & Pencarian, reranker, model jawaban dan eskalasi tiket. Reranker tersedia
melalui konfigurasi; profil integrasi terakhir diuji dengan reranker nonaktif. \\
18 & Demo alur nyata, hasil uji, keterbatasan dan penyerahan kode. \\
\bottomrule
\end{tabularx}

\topic{4. Status hasil dan bukti pemeriksaan}
\begin{tabularx}{\linewidth}{@{}p{39mm}Y@{}}
\toprule
Area & Bukti yang tersedia \\
\midrule
Refactoring & Cakupan awal serta temuan audit A1/A2 telah ditutup; dokumentasi
dan verifier disertakan. \\
Backend Laravel & 227 tes / 1403 assertion lulus pada audit. \\
Python AI & 228 tes lulus pada audit. \\
Pipeline dokumen/chat & 16 pemeriksaan integrasi nyata pada layanan sementara,
termasuk OCR, koreksi, embedding, JSON/SSE dan pembatalan. \\
Pemulihan cleanup & 9 pemeriksaan termasuk setup, untuk pemulihan setelah
gangguan Chroma/API. \\
Form dan UI & 18 pemeriksaan browser terkelompok per run; 210 pengukuran
kontrol cocok. 21/24 pasangan screenshot identik; tiga selisih raster kecil ditinjau. \\
\bottomrule
\end{tabularx}

{\small Hasil di atas berasal dari run yang didokumentasikan pada 25--26 Oktober
2026. Tes unit menggunakan dependency tiruan; uji integrasi memakai data
sintetis dan layanan sementara. Angka tersebut bukan hasil baru pada server instansi.}

\topic{5. Batas penggunaan yang perlu dipahami}
\begin{itemize}
\item Jawaban AI tetap perlu dinilai terhadap dokumen dan pertanyaan instansi;
pengujian sintetis tidak membuktikan akurasi umum.
\item Pengiriman kanal eksternal ke penerima nyata dan penerapan server instansi
belum dibuktikan oleh penutupan refactoring.
\item Dua migrasi pembatalan/outbox diuji pada database sementara; rollout
operasional memerlukan persiapan operator.
\end{itemize}

\newpage
\topic{6. Demo untuk pembimbing}
Siapkan layanan, akun demo dan dokumen contoh sebelum pertemuan. Alokasi
berikut adalah susunan presentasi sekitar 14 menit, bukan janji waktu inferensi.

\begin{tabularx}{\linewidth}{@{}p{17mm}p{37mm}Y@{}}
\toprule
Alokasi & Kegiatan & Bukti yang ditunjukkan \\
\midrule
2 menit & Masalah dan tujuan & Hubungan kebutuhan dinas dengan ruang lingkup proyek. \\
4 menit & Dokumen & Contoh unggahan, hasil ekstraksi, koreksi dan basis pengetahuan. \\
4 menit & Pertanyaan & Satu jawaban dari dokumen beserta sumber, lalu pertanyaan di luar cakupan. \\
2 menit & Tindak lanjut & Tiket internal atau direktori, status layanan dan monitoring. \\
2 menit & Penyerahan & Hasil uji, lokasi kode, batas penggunaan dan masukan pembimbing. \\
\bottomrule
\end{tabularx}

Gunakan data yang boleh ditampilkan. Jika memakai hasil ekstraksi yang sudah
diproses, jelaskan hal itu. Catat keluaran yang tidak sesuai sebagai hasil
evaluasi. Untuk demo yang mengubah data, gunakan lingkungan demo terpisah.

\topic{7. Paket penyerahan yang disarankan}
\begin{itemize}
\item Ringkasan ini, PPT yang telah dipresentasikan, dan PDF laporan akademik final.
\item Kode ketiga aplikasi, modul baru, lock file, migrasi, tes dan dokumentasi.
\item Contoh konfigurasi, kebutuhan layanan/model, dan panduan menjalankan.
\item Data demo yang boleh dibagikan, bukti uji dan daftar tindak lanjut.
\end{itemize}

Bekukan versi kode dan pastikan file baru ikut dalam paket. Pilih berkas yang
dibagikan; environment berisi kredensial, token, database operasional, log dan
referensi laporan pihak lain tidak menjadi isi paket kode. Akses operasional
diserahkan terpisah kepada pihak yang berwenang bila diperlukan.

\topic{8. Hal yang diminta dari pembimbing}
\begin{itemize}
\item Konfirmasi kesesuaian alur dengan kebutuhan unit.
\item Prioritas dokumen dan pertanyaan untuk evaluasi berikutnya.
\item Penetapan penerima kode, operator dan lingkungan uji penerapan.
\item Catatan hasil review dan tindak lanjut yang disepakati.
\end{itemize}

\vfill
{\small\color{muted}\textbf{Panduan pembacaan repository:} mulai dari
\texttt{SERAH\_TERIMA.md}. Bukti rinci tersedia pada dokumen status,
audit dan penutupan refactoring di folder \texttt{docs/}. Ringkasan ini
adalah draf untuk review, bukan pernyataan sudah diterima atau disetujui pembimbing.}
\end{document}
